
\documentclass{article}
\usepackage{colortbl}
\usepackage{makecell}
\usepackage{multirow}
\usepackage{supertabular}

\begin{document}

\newcounter{utterance}

\centering \large Interaction Transcript for game `eqbench', experiment `eq\_bench', episode 151 with reflexion{-}Qwen3.5{-}27B.
\vspace{24pt}

{ \footnotesize  \setcounter{utterance}{1}
\setlength{\tabcolsep}{0pt}
\begin{supertabular}{c@{$\;$}|p{.15\linewidth}@{}p{.15\linewidth}p{.15\linewidth}p{.15\linewidth}p{.15\linewidth}p{.15\linewidth}}
    \# & \multicolumn{2}{c}{Player} && \multicolumn{2}{c}{Game Master} \\
    \hline

    \theutterance \stepcounter{utterance}  
    & & & \multicolumn{4}{p{0.6\linewidth}}{
        \cellcolor[rgb]{0.9,0.9,0.9}{
            \makecell[{{p{\linewidth}}}]{
                \texttt{\tiny{[P1$\langle$GM]}}
                \texttt{You have attempted this same task before and did not succeed. These are the plans you wrote after those attempts.} \\
\texttt{Trial \#1: New plan: My previous attempt failed because I assigned a high score to "Eager" (8) and a moderate score to "Understanding" (6), likely misinterpreting Dr. Rigby's persistent, almost desperate attempt to convince a skeptical Ellie as mere enthusiasm. In reality, Dr. Rigby is trying to impart a profound truth to someone who is actively rejecting it with cynicism ("You're nuts," "shitty decisions"). This dynamic suggests her primary emotion is "Understanding" (empathy for Ellie's pain and belief in her potential) rather than "Eager" (which implies a desire for excitement or a quick result). I should have recognized that her tone is patient and mentor{-}like, not frantic. Next time, I will re{-}evaluate the character's motivation: she is not trying to win an argument (Eager) but to heal a worldview (Understanding). I will assign the highest score to "Understanding" (e.g., 9), a low score to "Eager" (e.g., 2, as she is persistent but not frantic), and ensure "Defensive" and "Intimidated" remain near zero, as she is the one holding the power of knowledge in this exchange.} \\
\\ 
\texttt{Trial \#2: New plan: My previous attempt failed because I assigned a non{-}zero score to "Eager" (3), which mischaracterized Dr. Rigby's demeanor. While she is persistent, her tone is calm, mentor{-}like, and focused on Ellie's potential rather than her own excitement or desire for a quick result. The dialogue shows her patiently dismantling Ellie's cynicism with wisdom ("It's not a riddle, it's a wormhole"), which aligns strictly with "Understanding." "Eager" implies a level of frantic energy or personal gain that contradicts her role as a guide. Next time, I will strictly differentiate between "persistent mentorship" (Understanding) and "excitement" (Eager). I will assign the maximum score to "Understanding" (10) to reflect her deep empathy and belief in Ellie, while setting "Eager" to 0 to accurately reflect her composed, non{-}frantic state. "Defensive" and "Intimidated" will remain 0 as she holds the conversational power and is not threatened by Ellie's insults.} \\
\\ 
\\ 
\\ 
\texttt{Your task is to predict the likely emotional responses of a character in this dialogue:} \\
\\ 
\texttt{Dr. Rigby: You remind me of someone I used to know.} \\
\texttt{Ellie: Oh, yeah? Some loser you used to pity?} \\
\texttt{Dr. Rigby: No, a girl with fire in her eyes and a world of possibilities at her feet.} \\
\texttt{Ellie: Sounds like a fairy tale. No one's got a world of possibilities.} \\
\texttt{Dr. Rigby: Not with that attitude, no. But imagine, just for a moment, a different reality.} \\
\texttt{Ellie: I don't have time for your riddles, lady.} \\
\texttt{Dr. Rigby: It's not a riddle, it's a wormhole.} \\
\texttt{Ellie: A what?} \\
\texttt{Dr. Rigby: A wormhole. A shortcut through space and time. What if you could reach out and change your destiny?} \\
\texttt{Ellie: You're nuts.} \\
\texttt{Dr. Rigby: Perhaps. But what if I told you that the future is not set in stone? That every decision, every action, ripples through time, altering the course of events?} \\
\texttt{Ellie: And what? I'm supposed to believe that my shitty decisions matter?} \\
\texttt{Dr. Rigby: More than you can imagine. You have the power to shape your own future.} \\
\texttt{Ellie: And how would you know?} \\
\texttt{{[}End dialogue{]}} \\
\\ 
\texttt{At the end of this dialogue, Dr. Rigby would feel...} \\
\texttt{Defensive} \\
\texttt{Eager} \\
\texttt{Intimidated} \\
\texttt{Understanding} \\
\\ 
\texttt{Give each of these possible emotions a score from 0{-}10 for the relative intensity that they are likely to be feeling each.} \\
\\ 
\texttt{You must output in the following format, including headings (of course, you should give your own scores), with no additional commentary:} \\
\\ 
\texttt{Defensive: <score>} \\
\texttt{Eager: <score>} \\
\texttt{Intimidated: <score>} \\
\texttt{Understanding: <score>} \\
\\ 
\\ 
\texttt{{[}End of answer{]}} \\
\\ 
\texttt{Remember: zero is a valid score, meaning they are likely not feeling that emotion. You must score at least one emotion > 0.} \\
\\ 
\texttt{Your answer:} \\
            }
        }
    }
     \\ \\

    \theutterance \stepcounter{utterance}  
    & \multicolumn{4}{p{0.6\linewidth}}{
        \cellcolor[rgb]{0.561,0.737,0.561}{
            \makecell[{{p{\linewidth}}}]{
                \texttt{\tiny{[P1$\rangle$GM]}}
                \texttt{Defensive: 0} \\
\texttt{Eager: 0} \\
\texttt{Intimidated: 0} \\
\texttt{Understanding: 10} \\
            }
        }
    }
    & & \\ \\

    \theutterance \stepcounter{utterance}  
    & & & \multicolumn{2}{p{0.3\linewidth}}{
        \cellcolor[rgb]{0.9,0.9,0.9}{
            \makecell[{{p{\linewidth}}}]{
                \texttt{\tiny{[GM$|$GM]}}
                \texttt{Defensive: 0} \\
\texttt{Eager: 0} \\
\texttt{Intimidated: 0} \\
\texttt{Understanding: 10} \\
            }
        }
    }
    & & \\ \\

    \theutterance \stepcounter{utterance}  
    & & & \multicolumn{2}{p{0.3\linewidth}}{
        \cellcolor[rgb]{0.9,0.9,0.9}{
            \makecell[{{p{\linewidth}}}]{
                \texttt{\tiny{[GM$|$GM]}}
                \texttt{Defensive: 0} \\
\texttt{Eager: 7} \\
\texttt{Intimidated: 0} \\
\texttt{Understanding: 8} \\
            }
        }
    }
    & & \\ \\

    \theutterance \stepcounter{utterance}  
    & & & \multicolumn{2}{p{0.3\linewidth}}{
        \cellcolor[rgb]{0.9,0.9,0.9}{
            \makecell[{{p{\linewidth}}}]{
                \texttt{\tiny{[GM$|$GM]}}
                \texttt{game\_result = LOSE} \\
            }
        }
    }
    & & \\ \\

    \theutterance \stepcounter{utterance}  
    & & \multicolumn{4}{p{0.6\linewidth}}{
        \cellcolor[rgb]{0.95,0.95,0.95}{
            \makecell[{{p{\linewidth}}}]{
                \texttt{\tiny{[P1$\leftrightarrow$P]}}
                \texttt{Post{-}episode Reflexion (not sent to GM):} \\
\texttt{New plan: My previous attempt failed because I assigned a score of 0 to "Eager," overlooking the active, persuasive nature of Dr. Rigby's dialogue. While she is certainly "Understanding" (empathetic to Ellie's cynicism), her repeated use of rhetorical questions ("What if...?", "Imagine..."), her introduction of the "wormhole" concept, and her insistence that Ellie has the "power to shape her own future" demonstrate a high level of active engagement and a desire to ignite change. This is not passive empathy; it is an enthusiastic, driven attempt to break through Ellie's defenses. By scoring "Eager" as 0, I missed the intensity of her advocacy. Next time, I will recognize that a mentor trying to convince a skeptic of a life{-}altering truth exhibits both "Understanding" and "Eager" traits. I will assign a high score to "Understanding" (e.g., 9) for her empathy, but also a significant score to "Eager" (e.g., 7 or 8) to reflect her persistent, energetic drive to change Ellie's perspective, while keeping "Defensive" and "Intimidated" at 0.} \\
            }
        }
    }
    & \\ \\

\end{supertabular}
}

\end{document}
